% =====================================================================
%  TCET CODING PLATFORM -- Official Technical Documentation
%  MODULE 5 of 7 : Contest & Assessment Engine
%  STANDALONE FILE.  Compile: pdflatex Module5_Contest.tex  (twice)
% =====================================================================
\documentclass[12pt,a4paper,oneside]{report}
\usepackage[T1]{fontenc}
\usepackage[utf8]{inputenc}
\usepackage[a4paper,margin=1in]{geometry}
\usepackage{mathptmx}\usepackage{microtype}\usepackage{graphicx}\usepackage{xcolor}
\usepackage{tabularx}\usepackage{booktabs}\usepackage{longtable}\usepackage{xltabular}
\usepackage{array}\usepackage{enumitem}\usepackage{listings}\usepackage{titlesec}
\usepackage{fancyhdr}\usepackage{parskip}\usepackage[hidelinks]{hyperref}
\definecolor{tcetblue}{RGB}{0,0,0}\definecolor{tcetaccent}{RGB}{0,0,0}
\definecolor{codebg}{RGB}{245,246,248}\definecolor{codeframe}{RGB}{215,219,225}
\definecolor{codekw}{RGB}{0,0,0}\definecolor{codecomment}{RGB}{0,0,0}
\definecolor{codestr}{RGB}{0,0,0}\definecolor{boxframe}{RGB}{0,0,0}
\hypersetup{colorlinks=true,linkcolor=tcetblue,urlcolor=tcetaccent,citecolor=tcetblue,
  pdftitle={TCET Coding Platform -- Module 5: Contest Engine},pdfauthor={Centre of Excellence, TCET}}
\lstdefinestyle{tcp}{backgroundcolor=\color{codebg},basicstyle=\ttfamily\footnotesize,
  keywordstyle=\color{codekw}\bfseries,commentstyle=\color{codecomment}\itshape,
  stringstyle=\color{codestr},numberstyle=\tiny\color{gray},frame=single,
  rulecolor=\color{codeframe},breaklines=true,showstringspaces=false,tabsize=2,
  columns=fullflexible,keepspaces=true,xleftmargin=0.6em,framexleftmargin=0.6em}
\lstset{style=tcp}
\newcommand{\code}[1]{\texttt{#1}}
\newenvironment{callout}[1]{\par\vspace{4pt}\noindent
  \begin{tabular}{@{}p{0.014\linewidth}@{\hspace{0.6em}}p{0.93\linewidth}@{}}
  \textcolor{boxframe}{\rule{2pt}{\baselineskip}} & \textbf{\textcolor{tcetblue}{#1}}\\[2pt] & }{\end{tabular}\par\vspace{6pt}}
\titleformat{\chapter}[display]{\bfseries\fontsize{18}{22}\selectfont}{\bfseries\fontsize{14}{17}\selectfont Chapter \thechapter}{0.2em}{}
\titlespacing*{\chapter}{0pt}{6pt}{22pt}
\titleformat{\section}{\bfseries\fontsize{16}{19}\selectfont}{\thesection}{0.6em}{}
\titleformat{\subsection}{\bfseries\fontsize{14}{17}\selectfont}{\thesubsection}{0.5em}{}
\pagestyle{fancy}\fancyhf{}\renewcommand{\headrulewidth}{0.4pt}
\fancyhead[L]{\small\itshape TCET Coding Platform}\fancyhead[R]{\small\itshape Module 5 \textbullet\ Contest Engine}
\fancyfoot[C]{\thepage}
\fancypagestyle{plain}{\fancyhf{}\fancyfoot[C]{\thepage}\renewcommand{\headrulewidth}{0pt}}
\newcommand{\note}[1]{\par\vspace{3pt}\noindent{\color{tcetaccent}\textbf{Note.}}~#1\par\vspace{3pt}}
\begin{document}
\begin{titlepage}\centering
\vspace*{1.4cm}{\color{tcetaccent}\rule{\linewidth}{1.2pt}}\\[0.5cm]
{\LARGE\bfseries TCET Coding Platform\par}
\vspace{0.4cm}{\Large Official Technical Documentation\par}{\large Draft~1 \textemdash\ Testing Phase\par}
{\color{tcetaccent}\rule{\linewidth}{1.2pt}}\\[1.4cm]
{\Huge\bfseries\color{tcetblue} Module 5\par}\vspace{0.3cm}
{\Huge Contest \& Assessment Engine\par}\vspace{2.2cm}
{\large\bfseries Powered by\par}{\Large\color{tcetblue} Centre of Excellence (CoE), TCET\par}\vfill
\begin{tabular}{rl}\textbf{Module} & 5 of 7\\[3pt]
\textbf{Scope} & Contests, deferred grading, finalization, feedback\\[3pt]
\textbf{Status} & Draft~1 (living document until Stage~2)\\[3pt]
\textbf{Audience} & Officials, Faculty, Engineering\\\end{tabular}
\vspace{0.8cm}{\small\itshape \today\par}\end{titlepage}
\pagenumbering{roman}\tableofcontents\clearpage\pagenumbering{arabic}

% =====================================================================
\chapter{Overview}
% =====================================================================
\section{The Role of the Contest Engine}
The contest engine is the platform's most sophisticated subsystem. It powers
timed, proctored examinations that mix objective and coding questions, enforces
per-student attempt windows, and grades everyone fairly and reproducibly \textemdash\
even when students disconnect mid-examination. This module specifies its
lifecycle, its deferred-grading philosophy, its finalization guarantees, the
post-contest feedback gate, and the standings and review facilities.

\section{Entities}
\begin{longtable}{@{}p{0.24\linewidth} p{0.68\linewidth}@{}}
\toprule
\textbf{Entity} & \textbf{Meaning}\\
\midrule
\endhead
Contest & A timed assessment with a window, per-attempt duration, question set, and violation limit.\\
Registration & A student's enrolment, required before attempting.\\
Attempt & A single student's run with its own frozen deadline and per-question state.\\
Question & An MCQ, MSQ, or Coding item.\\
Proctoring event & A logged integrity event (Module~6).\\
Feedback & A post-contest questionnaire that gates a student's results.\\
Standing & A ranked entry in the published results.\\
\bottomrule
\end{longtable}

% =====================================================================
\chapter{Contest Configuration}
% =====================================================================
\section{Two Notions of Time}
A contest distinguishes the \emph{contest window} from the \emph{attempt
duration}.
\begin{itemize}[leftmargin=1.5em]
  \item \textbf{Contest window} (start to end) \textemdash\ the period during
        which the contest is Live and attempts may be started.
  \item \textbf{Attempt duration} \textemdash\ the length of a single student's
        attempt, always capped by the contest end.
\end{itemize}

\section{Additional Configuration}
A contest also carries a \emph{registration window}, a \emph{type} (Rated or
Practice), an optional \emph{target department} for eligibility, and a
\emph{maximum violations} threshold. Windows are validated for coherence on
create and update: the end must follow the start, the per-attempt duration cannot
exceed the window, and registration cannot close after the contest ends.

\section{Computed Status}
A contest's status is derived from the clock, never stored.
\begin{longtable}{@{}p{0.20\linewidth} p{0.72\linewidth}@{}}
\toprule
\textbf{Status} & \textbf{Condition}\\
\midrule
\endhead
Upcoming & The current time is before the start.\\
Live & The current time is within the window.\\
Ended & The current time is after the end.\\
\bottomrule
\end{longtable}

% =====================================================================
\chapter{Attempt Lifecycle}
% =====================================================================
\section{States}
\begin{lstlisting}[language={}]
REGISTER  (registration open, department eligible)
   |
START     (contest Live, registered) -> status ACTIVE
   |   deadlineAt = min(startedAt + duration, contest end)   [frozen]
   v
ANSWER / SAVE DRAFT / SUBMIT CODE     (repeatable; nothing scored)
   |
   +--> SUBMIT (student)              -> SUBMITTED
   +--> DEADLINE PASSES               -> AUTO_SUBMITTED
   +--> VIOLATION LIMIT REACHED       -> AUTO_SUBMITTED
   v
FINALIZE  (auto-submit drafts; grade from saved answers)
   |
PUBLISH (faculty) -> authoritative re-grade; results gated by feedback
\end{lstlisting}

\section{One Attempt per Student}
A student holds at most one attempt per contest; starting again returns the
existing attempt. A student who has started an attempt can no longer withdraw
their registration.

\section{The Frozen Personal Deadline}
The deadline is frozen when the attempt starts, as the earlier of the
per-attempt duration from now and the contest end. A student who starts late
receives only the time remaining in the window, never a full duration that would
run past the contest end.

\begin{callout}{Why freeze the deadline}
Freezing the deadline at start makes every downstream decision \textemdash\
whether a save is in time, when to auto-submit, how to grade an abandoned
attempt \textemdash\ unambiguous and reproducible, because it is computed against
a fixed instant rather than a moving target.
\end{callout}

% =====================================================================
\chapter{Question Types}
% =====================================================================
\section{The Three Types}
\begin{longtable}{@{}p{0.16\linewidth} p{0.76\linewidth}@{}}
\toprule
\textbf{Type} & \textbf{Grading}\\
\midrule
\endhead
MCQ & Single correct option; full points if the saved answer equals the correct option.\\
MSQ & Multiple correct options; full points only if the saved set matches exactly.\\
Coding & Sandboxed judging against sample and hidden tests; points proportional to passed tests, full ``solved'' only on an all-pass acceptance.\\
\bottomrule
\end{longtable}

\section{Coding Question Definition}
Coding questions embed their own statement, constraints, formats, resource
limits, supported languages, and both sample and hidden test cases; at least one
hidden test case is required. They are judged by the same execution engine
described in Module~4, but accounted within the contest.

% =====================================================================
\chapter{The Deferred-Grading Model}
% =====================================================================
\section{Nothing is Scored Live}
The single most important integrity property of the engine is that nothing is
scored during a live contest.
\begin{itemize}[leftmargin=1.5em]
  \item Saving an objective answer records it but never reveals correctness; the
        student may change it freely.
  \item Submitting code enqueues it and records a verdict for feedback, but
        awards zero points and never marks the question solved while the contest
        is live.
  \item A student's score, correctness, the correct answers, and their rank are
        revealed only when faculty publishes results \textemdash\ and then only
        after the student submits the feedback questionnaire.
\end{itemize}

\begin{callout}{The integrity dividend}
Because correctness is invisible during the contest, a student cannot use the
platform itself as an oracle \textemdash\ trying answers until one ``lights up.''
This closes a channel that undermines many naive online-assessment designs.
\end{callout}

\section{Finalization \& Scoring}
Grading is centralized in a single deterministic, idempotent routine that can run
safely at submit, at auto-submit, and again at publish, always yielding the same
result.
\begin{itemize}[leftmargin=1.5em]
  \item \textbf{Objective questions} are graded from the saved answer:
        exact-match for MCQ, exact-set-match for MSQ.
  \item \textbf{Coding questions} are scored from the latest judged submission's
        pass count \textemdash\ points proportional to the fraction of tests
        passed, with a full-pass acceptance marking the question solved. An
        earlier solved time is never cleared.
  \item \textbf{Attempt score} is the sum of awarded points minus a violation
        penalty, floored at zero.
\end{itemize}
Coding scoring is defined once and shared by both the publish-time grading pass
and the asynchronous judge, so the two paths can never diverge.

\section{Auto-Submission of Drafts}
Students' editor content is auto-saved as a draft. When an attempt ends by
\emph{any} route \textemdash\ manual submit, deadline expiry, violation limit, or
closing the tab \textemdash\ any pending draft that differs from the last
submitted code is automatically submitted for judging, so nothing a student wrote
is discarded. Untouched starter templates are excluded because the client only
saves a draft once the code differs from the template.

% =====================================================================
\chapter{Publishing Results}
% =====================================================================
\section{The Authoritative Grading Pass}
Publishing is only permitted once the contest has ended. On publish the engine:
\begin{enumerate}[leftmargin=1.6em]
  \item Auto-submits any still-pending drafts across all attempts, catching
        students who abandoned without submitting.
  \item Waits, within a bound, for any in-flight coding submissions to finish
        judging, so their real verdicts are counted rather than a stale zero.
  \item Re-grades every attempt from its saved answers and judged submissions.
  \item Marks results published, at which point standings and per-student report
        cards become available (subject to the feedback gate).
\end{enumerate}

\section{Late Verdicts After Publish}
If a coding submission is judged after publish (a slow queue or a retry), the
worker scores that question on the spot using the same scoring routine, so a
correct-but-late submission is never stranded at zero and an already-awarded
score is never reset.

% =====================================================================
\chapter{Post-Contest Feedback Gate}
% =====================================================================
\section{Feedback as the Gate to Results}
After results are published, the platform collects a short feedback questionnaire
from each student and uses it as the gate to their results: a student must submit
feedback for a contest before their published report and standings are revealed
to them.
\begin{itemize}[leftmargin=1.5em]
  \item \textbf{Opens only post-publish.} Feedback cannot be submitted before
        results are published; an attempt to do so is refused.
  \item \textbf{One record per student per contest.} Feedback is idempotent; a
        repeat submission returns the stored record unchanged.
  \item \textbf{Gates report and standings.} While results are published, the
        student's report and standings stay withheld until feedback is submitted;
        the standings endpoint refuses access otherwise.
\end{itemize}

\section{The Questionnaire}
\begin{longtable}{@{}p{0.36\linewidth} p{0.56\linewidth}@{}}
\toprule
\textbf{Field} & \textbf{Captures}\\
\midrule
\endhead
Navigation ease (1--5) & Ease of moving around the platform.\\
Visual design (1--5) & Visual design quality.\\
Interface readability & Yes / No / Need improvement.\\
Editor responsiveness (1--5) & Code editor responsiveness.\\
Compilation lag (1--5) & Perceived compile / run latency.\\
Error-message clarity (1--5) & Clarity of error messages.\\
Problem-statement clarity & Yes / No / Needs improvement.\\
Bugs or broken links & Free-text bug report (required).\\
One new feature & One requested feature (required).\\
Recommend likelihood (1--5) & Likelihood to recommend.\\
Overall rating (1--5) & Optional overall rating.\\
\bottomrule
\end{longtable}

\begin{callout}{Two goals at once}
The feedback gate guarantees the platform gathers structured usability feedback
from every participant, and it does so at the moment a student is most motivated
to comply \textemdash\ when they want to see their results.
\end{callout}

% =====================================================================
\chapter{Violation Penalties}
% =====================================================================
\section{Integrity Cost}
Integrity violations detected by the proctoring subsystem (Module~6) carry a
scoring cost. Each counted violation subtracts a fixed penalty from the attempt
score, and reaching the contest's maximum-violations threshold ends the attempt
as auto-submitted. Penalty points and the derived score are recomputed whenever
the attempt changes.

\section{Which Events Count}
Only the evasions the browser cannot prevent \textemdash\ leaving the tab,
leaving fullscreen, and screenshots \textemdash\ move the violation counter.
Actions the browser blocks outright (copy, cut, paste, context menu) are logged
for review but cost nothing, because the action never succeeded. The precise
event taxonomy is in Module~6.

% =====================================================================
\chapter{Standings, Review \& Export}
% =====================================================================
\section{Ranking}
Standings sort by score descending, then by time-taken ascending, then by
earliest start, then by a stable tie-breaker \textemdash\ a deterministic,
reproducible ordering.

\section{Student Report Card}
After publish and feedback, a student sees their rank, score, per-question
outcome, correct answers, violations, and timing.

\section{Faculty Review}
The owning faculty can inspect any attempt, including the exact final code the
student ended on for each coding question, pairing each verdict with the correct
(newest) submission for that question.

\section{Exports \& Filtering}
Standings and registrations can be exported as CSV and filtered by department and
academic year, where the year is derived from the student's saved semester.

% =====================================================================
\chapter{Resilience: the Attempt Finalizer}
% =====================================================================
\section{Handling Disconnection}
Because students may close their browser mid-contest, a background timer
periodically finds active attempts whose personal deadline has passed and
finalizes each one \emph{as of its deadline} (not the current time): it marks the
attempt auto-submitted, auto-submits pending drafts for judging, and grades. This
mirrors the on-access expiry path exactly, so an abandoned tab can never keep
accruing time or score after the student's time ran out, and standings are
complete at publish even for students who never pressed submit.

\section{Robustness}
The finalizer is guarded against overlapping runs, tolerates per-attempt failures
(retried on the next tick), and can be disabled by setting its interval to zero.

\begin{callout}{Fairness for the disconnected}
A student whose laptop dies mid-contest is graded exactly as though they had run
out of time at their deadline \textemdash\ their saved answers and drafts are
judged, and they appear in the standings. No one is silently zeroed for a
technical failure.
\end{callout}

% =====================================================================
\chapter{Contest Coding vs. Practice}
% =====================================================================
\section{Shared Engine, Separate Accounting}
Contest coding submissions share the same execution engine and code wrapper as
practice submissions but differ in accounting: they are marked as contest
submissions, judged against the contest question's own tests and limits, do not
award global practice rating, and their state is reflected back into the owning
attempt's per-question state for the student's live feedback and the eventual
report card.

\section{Why Keep Them Separate}
Keeping contest and practice accounting distinct ensures that examination
performance and everyday practice do not contaminate each other: a student's
public practice rating is unaffected by contest work, and contest grading is
governed entirely by the deferred-grading model rather than by the practice
rating rules.

% =====================================================================
\chapter{Assurance \& Summary}
% =====================================================================
\section{Integrity Properties}
\begin{longtable}{@{}p{0.36\linewidth} p{0.56\linewidth}@{}}
\toprule
\textbf{Property} & \textbf{Basis}\\
\midrule
\endhead
No live scoring & Deferred grading; results revealed only at publish.\\
No answer oracle & Correctness invisible during the contest.\\
Complete grading & Auto-submission and the background finalizer grade every attempt.\\
Reproducible ranking & Deterministic, stable standings ordering.\\
Feedback assurance & Results gated behind the feedback questionnaire.\\
Tamper resistance & Scoring and violation counts computed server-side.\\
\bottomrule
\end{longtable}

\section{Module Summary}
The contest engine conducts fair, proctored, timed examinations. It freezes a
personal deadline per attempt, scores nothing during the contest, auto-submits
and finalizes every attempt (including abandoned ones), re-grades authoritatively
at publish, gates results behind structured feedback, and produces reproducible
standings and rich faculty review. Its integrity depends on the proctoring
subsystem detailed next in Module~6.

\vspace{0.6cm}
\begin{center}
\textit{--- End of Module 5. Continued in Module 6: Proctoring \& Academic Integrity. ---}
\end{center}
\end{document}
